\section*{अध्याय \ref{ch:b2:living-soil} --- जीवित मिट्टी}

\begin{solution}{exo:b2:living-soil:1}
O: कूड़ा तथा अधकचरा क्षयित कार्बनिक पदार्थ, जहाँ अपघटन
शुरू होता है। A: \omterm{def:b2:living-soil:soil}{ह्यूमस} से मिली खनिज \omterm{def:b2:living-soil:soil}{मिट्टी}, गहरी तथा कणिका-संरचना वाली,
जो अधिकांश जड़ें, जीव, पोषक तथा जल थामे है। B: जमाव का संस्तर,
जहाँ ऊपर से निक्षालित मृत्तिका, लोहे के ऑक्साइड तथा कार्बोनेट
जमा होते हैं। C: अपक्षयित जनक पदार्थ, किसी महीन आधात्री में
चट्टान-टुकड़े, जहाँ नई खनिज \omterm{def:b2:living-soil:soil}{मिट्टी} छूटती है। R:
बिना-अपक्षयित आधार-चट्टान।
\end{solution}

\begin{solution}{exo:b2:living-soil:2}
मृत्तिका तथा \omterm{def:b2:living-soil:soil}{ह्यूमस} की सतहों पर ऋण आवेश विनिमेय
धनायन थामे रखते हैं जिन्हें जड़ें प्रोटॉनों के बदले ले सकती हैं। बालू के कण
क्वार्ट्ज़ हैं, अनावेशित तथा मोटे, जिनकी प्रति ग्राम सतह थोड़ी है
और जिनमें कहने लायक \omterm{def:b2:living-soil:soil}{ह्यूमस} नहीं। चूना डालना कैल्सियम कार्बोनेट जोड़ता है: वह
कैल्सियम उन हाइड्रोजन तथा ऐलुमिनियम आयनों को हटा देता है जो किसी अम्लीय \omterm{def:b2:living-soil:soil}{मिट्टी} के
विनिमय-स्थल भर देते हैं, वह कार्बोनेट उस अम्लता को उदासीन करता है,
और वे स्थल किसी पोषक धनायन से फिर भर जाते हैं।
\end{solution}

\begin{solution}{exo:b2:living-soil:3}
अपघटक: जीवाणु (\emph{Bacillus}) तथा कवक
(\emph{Trichoderma}); सूक्ष्मजीव-चरने वाले: \omterm{def:b2:unicellular-diversity:protist}{प्रोटिस्ट} (अमीबा) तथा
जीवाणु-भक्षी सूत्रकृमि; अपरदभक्षी: केंचुए, सहस्रपाद,
संतुलक, वुडलाइस; परभक्षी: माइट, परभक्षी सूत्रकृमि,
कनखजूरे, ज़मीनी भृंग; और शीर्ष पर छछूंदर या छछूँदरी। वे
चरने वाले, जिनका $\mathrm{C:N}$ अपने जीवाणु-शिकार से ऊँचा है,
अतिरिक्त नाइट्रोजन अमोनियम के रूप में निकाल देते हैं: यानी वे
उस नाइट्रोजन को खनिजित करते हैं जिसे उन जीवाणुओं ने बंद कर रखा था।
\end{solution}

\begin{solution}{exo:b2:living-soil:4}
अर्बुस्कुली: वह कवक (कोई ग्लोमेरोमाइसीट) जड़ के
वल्कुट-कोशिकाओं में घुसता है और अर्बुस्कुल बनाता है; और अधिकांश शाकों तथा फ़सलों में और
उष्णकटिबंधीय वृक्षों में मिलती है। बाह्यकवकमूली: वह कवक (बेसिडियोमाइसीट तथा
ऐस्कोमाइसीट, जिनमें बहुत-से कुकुरमुत्ते हैं) जड़-सिरे को ग़िलाफ़ पहनाता है और
उसकी कोशिकाओं के बीच बिना उनमें घुसे बढ़ता है; और समशीतोष्ण
तथा बोरीय वृक्षों में मिलती है। दोनों में वह पौधा शर्करा देता है और वह कवक
फ़ॉस्फ़ेट, नाइट्रोजन, जल, सूक्ष्म-पोषक तथा रोगजनकों से कुछ
रक्षा।
\end{solution}

\begin{solution}{exo:b2:living-soil:5}
$L^{*} = I/k = 4/0.8 = \qty{5}{t/ha}$; \omterm{def:b2:biogeochemical-cycles:reservoir}{निवास-काल} $1/k = 1.25$
वर्ष; और $\ln 10/0.8 = 2.9$ वर्ष बाद \qty{90}{\%} हानि।
\end{solution}

\begin{solution}{exo:b2:living-soil:6}
खाए गए कार्बन के प्रति \qty{100}{kg} पर वे सूक्ष्मजीव \qty{40}{kg} रखते हैं और
उन्हें $40/8 = \qty{5}{kg}$ नाइट्रोजन चाहिए। पुआल $100/80 =
\qty{1.25}{kg}$ देता है: यानी कार्बन के प्रति \qty{100}{kg} पर \qty{3.75}{kg} का
स्थिरीकरण। क्लोवर $100/15 = \qty{6.7}{kg}$ देता है: यानी \qty{1.7}{kg} का
खनिजीकरण।
\end{solution}

\begin{solution}{exo:b2:living-soil:7}
वर्ष भर में कीड़ों से होकर \qty{20}{t/ha} \omterm{def:b2:living-soil:soil}{मिट्टी}। ऊपर के
\qty{20}{cm} का भार $0.2\times 10^{4}\times 1.3 = \qty{2600}{t/ha}$ है:
यानी एक बार गुज़रने को 130 वर्ष। डार्विन के प्रति एकड़ दस टन यानी
\qty{25}{t/ha} हैं, यानी वही कोटि --- उनका ``हर कुछ वर्षों में'' उस
सतही \omterm{def:b2:living-soil:soil}{मिट्टी} को कहता था, यानी कुछ सेंटीमीटर, पूरी ऊपरी \omterm{def:b2:living-soil:soil}{मिट्टी} को नहीं।
\end{solution}

\begin{solution}{exo:b2:living-soil:8}
जीवाणु: प्रति ग्राम $10^{9}\times 10^{-12} = \qty{1}{mg}$। तंतु:
आयतन $\pi(2\times 10^{-4})^{2}\times 10^{4} = \qty{1.26e-3}{cm^3}$,
यानी प्रति ग्राम द्रव्यमान \qty{1.4}{mg}। प्रति हेक्टेयर ($2.6\times 10^{9}$ g):
\qty{2.6}{t} जीवाणु तथा \qty{3.6}{t} कवक।
\end{solution}

\begin{solution}{exo:b2:living-soil:9}
$H^{*} = 1.5/0.02 = \qty{75}{t/ha}$ कार्बन। जोता गया: $1.5/0.04
= \qty{37.5}{t/ha}$; हानि \qty{37.5}{t/ha}, जिसका आधा
$\ln 2/0.04 = 17$ वर्ष बाद चला जाता है।
\end{solution}

\begin{solution}{exo:b2:living-soil:10}
फ़ॉस्फ़ेट: ऋतु भर में $\sqrt{2\times 10^{-13}\times 8.64\times 10^{6}} =
\qty{1.3}{mm}$ --- यानी अगली से \qty{1}{cm} दूर कोई जड़
उनके बीच की \omterm{def:b2:living-soil:soil}{मिट्टी} का आठवाँ भाग छूती है, जबकि \qty{0.1}{mm}
दूर तंतु उस सबको छूते हैं। नाइट्रेट: \qty{4}{cm}, यानी उस जड़-अंतराल से
दूर: वह स्वयं ही उस जड़ तक आ जाता है (और जल के साथ
द्रव्यमान-प्रवाह से भी), इसलिए कवकमूल नाइट्रेट के लिए थोड़ा ही जोड़ती हैं और
फ़ॉस्फ़ेट के लिए बहुत।
\end{solution}

\begin{solution}{exo:b2:living-soil:11}
द्रव्यमान $0.25\times 10^{4}\times 1.3 = \qty{3250}{t/ha}$। वर्ष भर में शुद्ध हानि
\qty{19}{t/ha}: यानी जीवनकाल 170 वर्ष। वर्ष भर में खोया कार्बन $20\times
0.03 = \qty{0.6}{t/ha}$। बिना-जुताई के नीचे शुद्ध हानि
\qty{1}{t/ha} है: यानी 3250 वर्ष।
\end{solution}

\begin{solution}{exo:b2:living-soil:12}
कूड़ा एक-दो वर्ष में फिरता है, \omterm{def:b2:living-soil:soil}{ह्यूमस} पचास में, वह खनिज A
संस्तर हज़ारों में (उसका बनना वर्ष भर में \qty{0.1}{mm} पर), और
लवण दशकों में जमा होता है तथा दशकों में बहा दिया जाता है यदि कोई
जल-निकास हो। कोई किसान कूड़े तथा फ़सल को ऋतु भर के भीतर अनुक्रिया करते
देखता है और \omterm{def:b2:living-soil:soil}{ह्यूमस} को अपने जीवन-भर के काम के भीतर (एक तिहाई का ह्रास
तीस वर्ष लेता है; और उबार उतना ही); जबकि उस खनिज \omterm{def:b2:living-soil:soil}{मिट्टी} का अपरदन
एक जीवनकाल में अदृश्य और दस में अप्रतिवर्ती है; और लवण
किसी पीढ़ी में दिखता है तथा, बिना जल-निकास के, स्थायी। जिन्हें कोई वार्षिक लेखा
अभिलिखित नहीं करता, वे ही \omterm{def:b2:living-soil:soil}{मिट्टी} के धीमे चर हैं।
\end{solution}

\begin{solution}{pb:b2:living-soil:1}
\textbf{1.} $0.25\times 10^{4}\times 1.3 = \qty{3250}{t/ha}$;
कार्बन $\qty{81}{t/ha}$।
\textbf{2.} $81/12 = \qty{6.8}{t/ha}$ नाइट्रोजन।
\textbf{3.} जीवाणु \qty{1}{mg/g}, $\qty{3.25}{t/ha}$; तंतु
$\pi(2\times 10^{-4})^{2}\times 2\times 10^{4}\times 1.1 =
\qty{2.8}{mg/g}$, \qty{9}{t/ha}: यानी सूक्ष्मजीवी जैवभार लगभग
\qty{12}{t/ha}।
\textbf{4.} लगभग \qty{14}{t/ha} जीवित पदार्थ (ताज़ा द्रव्यमान),
यानी उस ह्यूमस-कार्बन का कोई \qty{17}{\%} --- और उसका कुछ प्रतिशत भर
कार्बन के रूप में, क्योंकि जीव अधिकतर जल हैं।
\textbf{5.} सूक्ष्मजीवी कार्बन \qty{12}{t} का लगभग आधा, यानी \qty{6}{t};
वर्ष भर में दो बार फिरते और \qty{40}{\%} रखते हुए, वे
$2\times 6/0.4 = \qty{30}{t}$ कार्बन खपाते हैं और \qty{18}{t/ha}
कार्बन श्वसित करते हैं, यानी \qty{66}{t} $\mathrm{CO_2}$ (कोई मोटा आकलन: वह
धारण तथा फेर कच्चे हैं)।
\textbf{6.} उस फ़सल के \qty{5}{t} शुद्ध प्राथमिक उत्पादन से
तीन गुना: यानी वह आकलन बहुत ऊँचा है, क्योंकि वे मान्यताएँ (सारा
जैवभार दो बार फिरता, कुछ भी सुप्त नहीं) उदार हैं;
किसी फ़सल के नीचे असली \omterm{prop:b2:living-soil:humus}{मिट्टी-श्वसन} उस फ़सल के अपने
उत्पादन की कोटि का होता है, यानी प्रति हेक्टेयर प्रति वर्ष कई टन कार्बन। वह
\omterm{def:b2:living-soil:soil}{मिट्टी} अपने ऊपर के पौधों जितनी ही साँस लेती है।
\textbf{7.} कार्बन $5\times 0.45 = \qty{2.25}{t}$; नाइट्रोजन $5\times
0.03 = \qty{150}{kg}$; $\mathrm{C:N} = 15$।
\textbf{8.} खनिजित करता है: उन सूक्ष्मजीवों को $2250\times 0.4/8 =
\qty{112}{kg}$ नाइट्रोजन चाहिए और वह अवशेष 150 थामे है, इसलिए
\qty{38}{kg/ha} शुद्ध छूटती है।
\textbf{9.} 3 महीनों बाद $1 - e^{-0.5} = \qty{39}{\%}$; वर्ष भर बाद $1 - e^{-2}
= \qty{86}{\%}$।
\textbf{10.} $0.39\times 38 = \qty{15}{kg/ha}$।
\textbf{11.} उस फ़सल को \qty{150}{kg} चाहिए, और उसका अधिकांश अपने पहले
दो महीनों में; जबकि वह अवशेष तीन महीनों में \qty{15}{kg} और
कुल मिलाकर \qty{38}{kg} देता है --- यानी कोई पूरक, कोई आपूर्ति नहीं। उसका मूल्य
उस नाइट्रोजन में भी है जो उस क्लोवर ने स्थिर की और जिसे वह उस अवशेष के
अपने \qty{150}{kg} में थामे है, और जिसका अधिकांश \omterm{def:b2:living-soil:soil}{ह्यूमस} में घुसकर
वर्षों भर लौटता है।
\textbf{12.} पुआल: कार्बन \qty{2250}{kg}, नाइट्रोजन $2250/80 =
\qty{28}{kg}$; उन सूक्ष्मजीवों को \qty{112}{kg} चाहिए, इसलिए \omterm{def:b2:living-soil:soil}{मिट्टी} से
\qty{84}{kg/ha} स्थिरीकृत हो जाती है --- यानी उस फ़सल की ज़रूरत से आधी से अधिक।
उस किसान को उस पुआल के साथ नाइट्रोजन डालनी चाहिए, या उसे बोने से
अच्छी तरह पहले जोत देना चाहिए, या उसे सतह पर छोड़ देना चाहिए जहाँ वह धीरे क्षय होता है।
\textbf{13.} जोता गया: $1.2/0.025 = \qty{48}{t/ha}$; बिना-जुताई:
$1.2/0.015 = \qty{80}{t/ha}$।
\textbf{14.} $H(t) = 80 - 32\,e^{-0.015t}$: यानी 10 वर्षों बाद \qty{4.5}{t/ha}
का लाभ, और 50 बाद \qty{17}{t/ha}।
\textbf{15.} पहले वर्ष में $\mathrm{d}H/\mathrm{d}t = 0.015\times 32 =
\qty{0.48}{t/ha}$ कार्बन, यानी \qty{1.8}{t}
$\mathrm{CO_2}$।
\textbf{16.} पहले वर्ष में $0.48\times 1.5\times 10^{9} = \qty{0.7}{GtC}$,
यानी जीवाश्म उत्सर्जनों का \qty{7}{\%}; और पचासवें वर्ष तक
वह दर गिरकर $0.48\,e^{-0.75} = \qty{0.23}{t/ha}$ हो जाती है,
यानी \qty{0.34}{GtC}।
\textbf{17.} वह लाभ किसी नई \omterm{def:b2:biogeochemical-cycles:reservoir}{स्थायी अवस्था} तक की पहुँच है: वह
शून्य तक क्षीण होता है ज्यों ही $H$ \qty{80}{t/ha} तक पहुँचे, और यदि वह खेत फिर जोता जाए तो
पूरा लाभ दशकों के भीतर वायु को लौट जाता है
--- यानी कोई मिट्टी-पात्र कोई एक-बार का, प्रतिवर्ती संचरण है, न कि चलते उत्सर्जनों की कोई
स्थायी भरपाई।
\textbf{18.} $k_h = 0.02$: $H^{*} = \qty{60}{t/ha}$, यानी
\qty{20}{t/ha} की हानि, जो पहले पचास वर्षों के बिना-जुताई लाभ से भी अधिक है:
यानी तापन उस प्रबंधन को
उलट सकता है।
\textbf{19.} जोता गया: वर्ष भर में $15 - 0.5 = \qty{14.5}{t/ha}$; बिना-जुताई:
\qty{1}{t/ha}।
\textbf{20.} $3250/14.5 = 220$ वर्ष; $3250/1 = 3250$ वर्ष।
\textbf{21.} कार्बन $15\times 0.025 = \qty{0.38}{t/ha}$; नाइट्रोजन
वर्ष भर में $0.38/12 = \qty{31}{kg/ha}$।
\textbf{22.} जोते की \omterm{def:b2:biogeochemical-cycles:reservoir}{स्थायी अवस्था} पर ह्यूमस-क्षय $k_hH^{*} =
I = \qty{1.2}{t/ha}$ प्रति वर्ष है; और अपरदन उसका एक तिहाई और जोड़ देता है,
पर क्षय के उलट वह किसी निवेश से संतुलित नहीं --- वह उस
\omterm{def:b2:biogeochemical-cycles:reservoir}{भंडार} को बहा देता है।
\textbf{23.} आंशिक रूप से। किसी अवसाद या किसी जलाशय में गड़ा
अपरदित कार्बन लंबे समय तक क्षय से बचा रह सकता है (यानी कोई
दफ़न-पात्र); पर वे समुच्चय ढोने में टूट जाते हैं, उस कार्बन का अधिकांश
रास्ते में ऑक्सीकृत हो जाता है, और उस खेत की खोई उर्वरता
ऐसे उर्वरक से बदली जाती है जिसका निर्माण कार्बन उत्सर्जित करता है। अपरदन के कार्बन-प्रभाव का
वैश्विक चिह्न अब भी विवादित है।
\textbf{24.} कूड़ा $1/k = 0.5$ वर्ष; \omterm{def:b2:living-soil:soil}{ह्यूमस} हल के नीचे $1/k_h = 40$ वर्ष
और बिना-जुताई के नीचे 67; खनिज \omterm{def:b2:living-soil:soil}{मिट्टी} बनने के सापेक्ष $3250/0.5 =
6500$ वर्ष, और हल के नीचे हानि के सापेक्ष 220
वर्ष।
\textbf{25.} A संस्तर में कार्बन-भंडार \qty{81}{t/ha}; जीवनकाल
हल के नीचे 220 वर्ष, बिना-जुताई के नीचे 3250; और वह क्लोवर-अवशेष
लगभग \qty{38}{kg/ha} खनिज नाइट्रोजन देता है, जिनमें से 15
पहले तीन महीनों में।
\end{solution}
